% Wyświetlamy oba indeksy jeden po drugim
\chapter*{Indeks j}
\section*{Test}
\begin{multicols}{2}
\textbf{F}\nopagebreak

\begin{minipage}{\linewidth}
fizyka \newline
\hspace*{1.5em} grawitacja, \pageref{latex_test1} \newline
\hspace*{1.5em} grawitacja, \pageref{latex_test3} \newline
\end{minipage}
\begin{minipage}{\linewidth}
fizyka2 \newline
\hspace*{1.5em} mechanika kwantowa, \pageref{latex_test2} \newline
\hspace*{1.5em} mechanika kwantowa, \pageref{latex_test4} \newline
\end{minipage}

\end{multicols}

\section*{Test2}
\begin{multicols}{2}
\textbf{F}\nopagebreak

\begin{minipage}{\linewidth}
fizyka \newline
\hspace*{1.5em} grawitacja \dotfill \pageref{latex_test1} \newline
\hspace*{1.5em} grawitacja \dotfill \pageref{latex_test3} \newline
\end{minipage}
\begin{minipage}{\linewidth}
fizyka2 \newline
\hspace*{1.5em} mechanika kwantowa \dotfill \pageref{latex_test2} \newline
\hspace*{1.5em} mechanika kwantowa \dotfill \pageref{latex_test4} \newline
\end{minipage}
\end{multicols}

%%%%%

\markboth{\textbf{} (ゞ)}{\textbf{} (ゞ)}
\phantomsection
W fizyce klasycznej \textbf{grawitacja}\label{latex_test1} odgrywa kluczową rolę.
Teorię tę w wielkim stopniu rozwinął \textbf{Isaac Newton}\index[nazwiska]{Newton, Isaac}.

\phantomsection
Z kolei w skali mikro rządzi \textbf{mechanika kwantowa}\label{latex_test2},
której podwaliny kładł m.in. \textbf{Max Planck}\index[nazwiska]{Planck, Max}.

\newpage

\phantomsection
\markboth{\textbf{} (ゞ222)}{\textbf{} (ゞ222)}
222 W fizyce klasycznej \textbf{grawitacja}\label{latex_test3} odgrywa kluczową rolę.
Teorię tę w wielkim stopniu rozwinął \textbf{Isaac Newton}\index[nazwiska]{Newton, Isaac}.

\phantomsection
222 Z kolei w skali mikro rządzi \textbf{mechanika kwantowa}\label{latex_test4},
której podwaliny kładł m.in. \textbf{Max Planck}\index[nazwiska]{Planck, Max}.

%%%%%%%%%%%

